\section{Method}
\label{sec:method}

This section formalizes the \fastgrnn{} cell, the three-stage
compression pipeline applied on top of it, and the LUT-based
activation recipe that makes the result deployable on a
multiplier-less MCU. Fig.~\ref{fig:pipeline} summarizes the full
flow.

\subsection{\fastgrnn{} Cell Formulation}
\label{sec:method:cell}

Given an input $\mathbf{x}_t \in \mathbb{R}^{d}$ at time $t$ and a
previous hidden state $\mathbf{h}_{t-1} \in \mathbb{R}^{H}$, the
\fastgrnn{} cell~\cite{kusupati2018fastgrnn} computes a single gate
$\mathbf{z}_t$, a candidate update $\tilde{\mathbf{h}}_t$, and a
two-scalar interpolation $\mathbf{h}_t$:
\begin{align}
\mathbf{z}_t &= \sigma(\mathbf{W} \mathbf{x}_t + \mathbf{U} \mathbf{h}_{t-1} + \mathbf{b}_z) \label{eq:z}\\
\tilde{\mathbf{h}}_t &= \tanh(\mathbf{W} \mathbf{x}_t + \mathbf{U} \mathbf{h}_{t-1} + \mathbf{b}_h) \label{eq:htilde}\\
\mathbf{h}_t &= \bigl(\zeta(\mathbf{1} - \mathbf{z}_t) + \nu\bigr)\, \tilde{\mathbf{h}}_t + \mathbf{z}_t \odot \mathbf{h}_{t-1} \label{eq:hupdate}
\end{align}
where $\odot$ denotes element-wise multiplication and
$\zeta, \nu \in (0,1)$ are \emph{learned scalars} that interpolate
between leaky-integrator and standard gated update dynamics. The
two-scalar gate is the defining feature of \fastgrnn{}: it permits
LSTM-comparable expressiveness with a single weight pair
$(\mathbf{W}, \mathbf{U})$ shared between the gate and the candidate
update, in contrast to the three or four weight pairs of GRU or
LSTM. In our HAR setup, $d{=}3$ (tri-axial acceleration), $H{=}16$,
and the input sequence length is $T{=}128$ samples (one 2.56~s
window at 50~Hz). The choice $H{=}16$ is empirically justified in
Table~\ref{tab:hsize}: doubling the hidden size to $H{=}32$
triples the parameter count yet produces a \emph{lower} test F1,
and the larger model degrades further with longer training~--- the
canonical signature of overfitting on a dataset whose intrinsic
complexity the smaller cell already saturates.

\begin{table}[h]
\centering
\caption{Hidden-size selection on \hapt{}, full-rank \fastgrnn{}
(no low-rank, no sparsification). $H{=}32$ has roughly three times
the parameters of $H{=}16$ but is worse on every matched
configuration, and \emph{further} degrades from epoch~30 to~50,
indicating overfitting. $H{=}16$ continues to improve through
epoch~120 (val-selected ceiling F1~$=$~0.912; see
Section~\ref{sec:results:lsq}). $H{=}32$ was not extended to
120~epochs --- its monotone degradation over 30~$\rightarrow$~50
makes further training extremely unlikely to recover the gap.
We therefore fix $H{=}16$ for all downstream experiments.}
\label{tab:hsize}
\begin{tabular}{rrrrr}
\toprule
$H$ & Epochs & F1 & Accuracy & Params \\
\midrule
\textbf{16} & 30  & 0.847          & 0.854          & \textbf{440}   \\
\textbf{16} & 50  & 0.861          & 0.869          & \textbf{440}   \\
\textbf{16} & 120 & \textbf{0.912} & \textbf{0.913} & \textbf{440}   \\
32          & 30  & 0.833          & 0.845          & 1{,}384         \\
32          & 50  & 0.830          & 0.834          & 1{,}384         \\
32          & 120 & ---            & ---            & 1{,}384         \\
\bottomrule
\end{tabular}
\end{table}

Because the same weight pair $(\mathbf{W}, \mathbf{U})$ is reused in
both~\eqref{eq:z} and~\eqref{eq:htilde}, the unconstrained
parameter count of the cell is
\begin{align}
\#\text{params}
 &= \underbrace{Hd}_{\mathbf{W}} +
    \underbrace{H^2}_{\mathbf{U}} +
    \underbrace{2H}_{\mathbf{b}_z,\mathbf{b}_h} +
    \underbrace{2}_{\zeta,\nu} \nonumber\\
 &= 48 + 256 + 32 + 2 = 338.
\end{align}
A standard LSTM at the same hidden size requires four such
weight matrices ($\approx$\,1{,}280 parameters), so \fastgrnn{}
starts from a $\sim$$4\times$ advantage at the architectural level.
The low-rank and sparsification stages below reduce this further
by another order of magnitude.

\subsection{Low-Rank Weight Factorization}
\label{sec:method:lowrank}

The input matrix
$\mathbf{W} \in \mathbb{R}^{H \times d}$ and recurrent matrix
$\mathbf{U} \in \mathbb{R}^{H \times H}$ are each factored as a
product of two thin matrices:
\begin{align}
\mathbf{W} &= \mathbf{W}_1 \mathbf{W}_2^\top, \quad \mathbf{W}_1 \in \mathbb{R}^{H \times r_w}, \mathbf{W}_2 \in \mathbb{R}^{d \times r_w}, \\
\mathbf{U} &= \mathbf{U}_1 \mathbf{U}_2^\top, \quad \mathbf{U}_1 \in \mathbb{R}^{H \times r_u}, \mathbf{U}_2 \in \mathbb{R}^{H \times r_u}.
\end{align}
The factors $\mathbf{W}_1, \mathbf{W}_2, \mathbf{U}_1, \mathbf{U}_2$
are learned end-to-end. The recurrent matrix-vector product is
evaluated as $\mathbf{U}\mathbf{h} = \mathbf{U}_1(\mathbf{U}_2^\top \mathbf{h})$,
which costs $2 H r_u$ multiply-adds rather than $H^2$. After
sweeping $r_u \in \{4, 6, 8, 12\}$ across five random seeds each
(Section~\ref{sec:results:lowrank}), we select $r_w = 2, r_u = 8$.

\subsection{Iterative Hard Thresholding}
\label{sec:method:iht}

We sparsify all four factor matrices via iterative hard
thresholding~\cite{blumensath2009iht}: at each training step we
retain the top-$k$ magnitude entries of every weight tensor and
zero the rest. The target sparsity $s$ follows the cubic schedule
\begin{equation}
s_e = s \cdot \min\!\Bigl(1, \tfrac{e}{e_{\text{ramp}}}\Bigr)^{3}
\end{equation}
over training epochs $e$, with $e_{\text{ramp}} = 50$ followed by
50 epochs of mask-frozen fine-tuning. After sweeping
$s \in \{0.3, 0.5, 0.7, 0.9\}$ across five seeds
(Section~\ref{sec:results:sparsity}), we select $s = 0.5$
(283 nonzero parameters), a U-curve optimum that balances
accuracy against compressibility.

\subsection{Per-Tensor Q15 Quantization with Activation Calibration}
\label{sec:method:q15}

Each weight tensor $\mathbf{W}^{(\ell)}$ (where
$\ell \in \{W_1, W_2, U_1, U_2\}$) is quantized via
\begin{equation}
\hat{w}^{(\ell)}_{ij}
  = \mathrm{clip}\!\Bigl(
        \mathrm{round}\bigl(w^{(\ell)}_{ij} / s_\ell\bigr),
        -2^{15}, 2^{15}-1
    \Bigr),
\end{equation}
with a tensor-specific scale $s_\ell$ chosen so that
$\max_{ij}|w_{ij}| / s_\ell$ lies just below the Q15 ceiling.
At inference, the dequantized value $\hat{w}^{(\ell)}_{ij} \cdot s_\ell$
is used in the forward pass.

A naive extension of this scheme to activations
\emph{also} placed in Q15 $[-1, 1)$ proves catastrophic: the hidden
state $\mathbf{h}_t$ in our trained model attains magnitudes of
$\sim$62, an order of magnitude beyond the Q15 representable range.
Without intervention, F1 collapses from 0.918 to 0.16 and the
\textsc{standing} class disappears entirely
(Section~\ref{sec:results:quant}).

A straightforward fix is to switch to a fixed Q9.6 fixed-point
format, whose $\pm 64$ representable range covers the observed
$\mathbf{h}_t$ magnitudes by construction. This trades a uniform
safety margin for accuracy headroom on \emph{every} tensor,
regardless of its actual dynamic range. We instead use
\emph{per-activation calibration}: a calibration pass runs five
mini-batches of training data through the FP32 model, records the
empirical maximum of every intermediate tensor, applies a 10\%
headroom, and assigns each activation its own scale. The
per-tensor scheme generalizes Q9.6 -- it approaches the $\pm 64$
range adaptively for tensors that need it (e.g.\ $\mathbf{h}_t$)
while preserving the full Q15 resolution on tensors that do not
(e.g.\ the gate output $\mathbf{z}_t \in [0,1]$). With calibration, prediction agreement
between the integer C implementation and the FP32 reference is
100\% across the 3{,}399-window test set, and the C-side
macro F1 reaches 0.9176 -- the deployed-model number reported
throughout this paper.

\subsection{LUT-Based Activations}
\label{sec:method:lut}

The remaining cost on a multiplier-less target is the
$\sigma$ and $\tanh$ evaluations: each requires several
software-emulated transcendentals per call, and the cell evaluates
$2H$ such activations per sample ($2 \cdot 16 = 32$ at $H{=}16$,
times 128 samples per window = 4{,}096 activation calls per window).
We replace the runtime calls with a 256-entry lookup table
over the input domain $[-8, +8]$, stored in Flash:
\begin{equation}
\mathrm{LUT}[k]
  = f\!\Bigl( -8 + k \cdot \tfrac{16}{255} \Bigr), \quad
  k \in \{0, 1, \ldots, 255\},
\end{equation}
where $f \in \{\sigma, \tanh\}$. Inputs outside $[-8, +8]$
saturate to $\pm 1$, which is exact to floating-point precision
for both functions in those tails. Inside the domain, a single
linear interpolation between adjacent entries replaces the original
transcendental, reducing each activation to one comparison, two
indexed loads, a subtract, and a multiply-add. The two tables
together occupy 1~KB of Flash; the result is a
\textbf{30.5$\times$} end-to-end speedup on the \msp{}
(Section~\ref{sec:results:streaming}).

\begin{figure}[t]
\centering
\begin{tikzpicture}[
  node distance=4mm and 4mm,
  every node/.style={font=\scriptsize},
  stage/.style={draw, rounded corners=2pt, align=center,
                 minimum height=8mm, minimum width=14mm, fill=blue!8,
                 inner sep=2pt},
  deploy/.style={draw, rounded corners=2pt, align=center,
                 minimum height=8mm, minimum width=16mm, fill=green!10,
                 inner sep=2pt},
  arr/.style={-{Stealth[length=4pt]}, semithick},
]
\node[stage] (train)  {Float\\train};
\node[stage, right=of train]    (lr)    {Low-rank\\$r_w{=}2$, $r_u{=}8$};
\node[stage, right=of lr]       (sp)    {IHT\\sparsity 0.5};
\node[stage, right=of sp]       (q15)   {Q15\\+ calib.};

\node[deploy, below=6mm of lr]  (arduino) {Arduino Uno\\(AVR, $\times$HW)};
\node[deploy, below=6mm of q15] (msp)     {MSP430G2553\\(no $\times$HW)};

\node[align=center, below=2mm of arduino, font=\scriptsize\itshape, text=gray]
   (note1) {portable C + LUT};

\draw[arr] (train) -- (lr);
\draw[arr] (lr)    -- (sp);
\draw[arr] (sp)    -- (q15);
\draw[arr] (q15.south) -- ++(0,-3mm) -| (arduino.north);
\draw[arr] (q15.south) -- ++(0,-3mm) -| (msp.north);
\end{tikzpicture}
\caption{End-to-end \fastgrnn{} compression and deployment pipeline.
Float training $\rightarrow$ low-rank $\rightarrow$ sparse $\rightarrow$
calibrated Q15 weights $\rightarrow$ portable C with LUT activations
$\rightarrow$ Arduino and MSP430 binaries.}
\label{fig:pipeline}
\end{figure}

\subsection{Implementation}
\label{sec:method:impl}

Training is implemented in PyTorch~2.x; the deployed C inference
engine (\texttt{fastgrnn.cpp}) is a single $\sim$200-line
translation unit that compiles unmodified under both
\texttt{avr-gcc} (\arduino{}) and the Code Composer Studio
\texttt{msp430-elf-gcc} (\msp{}). All weights are stored in Flash
as a \texttt{const int16\_t} array, alongside the per-tensor
scales and the two activation LUTs. The runtime working set
(hidden state $\mathbf{h}$, gate $\mathbf{z}$, candidate
$\tilde{\mathbf{h}}$, output logits, and scratch) totals
$\sim$300~bytes, fitting comfortably within the 512~B SRAM budget
of the \msp{}.
